\clearpage
\fancyhead[L]{\small HARDING. HIMMELSATLAS.}
\SourceLeaf{515}
\begin{center}\small Göttingische gelehrte Anzeigen. Stück 61. Seite 601..607. 1809 April 17.\\[3pt]\rule{.72\linewidth}{.3pt}\end{center}

\textsc{K. L. Harding}. \emph{Neuer Himmelsatlas}. \emph{Erste Lieferung}, enthaltend die Blätter I. II. V und IX. Bei dem Verfasser und in Commission bei \textsc{Friedrich Perthes} in Hamburg, 1809.

Mit Vergnügen zeigen wir die Erstlinge von einer Unternehmung an, deren Ausführung Kenner und Liebhaber der Astronomie seit mehreren Jahren mit Verlangen entgegen sahen, welcher der Druck der Zeiten zwar schwere Hindernisse in den Weg legte, aber doch glücklicher Weise ihren Fortgang nicht gehemmt hat. Das Bedürfniss möglichst genauer und sehr detaillirter Himmelskarten, welches zwar schon früher, besonders bei der Erscheinung von Cometen, oft genug gefühlt wurde, ist seit Entdeckung der Ceres und der drei andern verschwisterten neuen Planeten für Jeden, der diese so merkwürdigen Weltkörper beobachten oder auch nur sehen will, so dringend nothwendig geworden, dass jedem Freunde der Sternkunde die Ausfüllung dieser Lücke, und eine solche Ausfüllung, wie diese vier ersten Blätter beurkunden, nicht anders als höchst willkommen sein kann. Den Entschluss zu dieser Unternehmung hatte unser verdienstvoller College bereits vor mehr als fünf Jahren, bald nach Entdeckung der Ceres und Pallas, gefasst: anfangs freilich nach einem etwas beschränkteren Plane, nach welchem nur diejenigen Zonen des Himmels bearbeitet werden sollten, in denen jene beiden Planeten uns erscheinen können: die Absteckung dieser Zonen gab damals zu einem interessanten analytischen Problem Veranlassung, welches Hr. Prof. \textsc{Gauss} in der Monatlichen Correspondenz (1804, August [S. 106 d. B.]) aufgelöst hat. Hr. Prof. \textsc{Harding} fing seine Arbeit unter den glücklichsten Auspicien an; die Entdeckung eines dritten neuen Planeten ({\LatitudeSymbols οὐ τύχης ἀλλ' ἀρετῆς εὐτυχούσης ἔργον}, wie \textsc{Delambre} mit Plutarchs Worten sagt) lohnte gleichsam im Voraus die ausgedehnte Arbeit, welcher er sich unterzog. Aber eben diese Entdeckung, und die 1807 dem Doctor \textsc{Olbers} geglückte Auffindung eines vierten Planeten, machten schon die Erweiterung des Plans nothwendig, und die Ueberzeugung von der Wichtigkeit der Sache, verbunden mit den vielfältig von andern Astronomen geäusserten Wünschen, bestimmten den Hrn. Prof. \textsc{Harding}, jene ursprünglichen Beschränkungen ganz fallen zu lassen und die Unternehmung auf den ganzen Fixsternhimmel auszudehnen, so weit er den Europäischen Beobachter interessirt, und so weit adäquate Hülfsmittel zu seiner Bearbeitung vorhanden sind. Von jenem ersten eingeschränkten Plane sind daher weiter keine Spuren zurückgeblieben, als in der Numerirung und Disposition der Karten, so wie es denn natürlich auch denen, welche sich für diese neue Darstellung des Sternhimmels interessiren, am gelegensten sein muss, dass diejenigen Blätter zuerst erscheinen, auf denen Stücke des Zodiakus der neuen Planeten vorkommen, um dem nothwendigsten Bedürfnisse so früh als möglich abzuhelfen.

Aufgenommen hat Hr. Prof. \textsc{Harding} in seine Karten \emph{alle} Sterne, von denen \emph{zuverlässige} Bestimmungen oder Beobachtungen vorhanden sind; eine beträchtliche Anzahl von Sternen, die noch in keinem Verzeichnisse vorkommen, hat er überdies noch nach eigenen Beobachtungen eingetragen, welche Bestimmungen er demnächst diesem Atlas beifügen wird. Wenn man erwägt, wie sehr viel für die bessere Kenntniss des Fixsternhimmels, besonders in den letzten Zeiten, geschehen ist, so wird man schon von selbst beurtheilen können, wie sehr dieser neue Atlas an Reichhaltigkeit alle seine Vorgänger übertreffen muss. So hat z. B. das Sternbild des Wassermanns in den \textsc{Flamsteed}'schen Himmelskarten 108 Sterne, in den grossen \textsc{Bode}'schen Himmelskarten 343, während wir auf Hrn. Prof. \textsc{Harding}'s Karten, Blatt IX und VIII (welches letztere zwar erst zur zweiten Lieferung bestimmt, aber jetzt auch bereits gestochen in unsern Händen ist) 1444 Sterne zählen. Es ist schwer zu entscheiden, ob dieser Reichthum selbst, oder die beispiellose Gewissenhaftigkeit, womit derselbe benutzt ist, für den grösse\SourceLeaf{516}ren Vorzug dieser Karten zu halten ist. Alle bisherigen Bearbeiter von Himmelskarten haben immer die Sterne so eingetragen, wie diejenigen Quellen, welche man für die besten hielt, sie angeben, und wenn sie auch, wie Hr. Prof. \textsc{Bode} bei seinem Atlas, durch sorgfältige Vergleichungen, Conjecturen und scharfsinnige Kritik eine Menge von Fehlern aller Art, welche in den Fixsternverzeichnissen bei einer so ungeheuren Menge Zahlen unvermeidlich waren, glücklich wegräumten: so konnte es doch nicht fehlen, dass sehr viele Unrichtigkeiten, zu deren Aufspürung jene Mittel nicht hinreichten, unbemerkt blieben, und in die Karten mit übergingen. Daher die vielen Sterne in den bisherigen Himmelskarten, die gar nicht am Himmel stehen und nie gestanden haben, dergleichen z. B. Hr. Prof. \textsc{Harding} unter den 343 \textsc{Bode}'schen Sternen des Wassermanns nicht weniger als 15, und in manchen andern Sternbildern verhältnissmässig noch weit mehr bemerkt hat. Das einzige Mittel, diesen Fehlern zu entgehen, bestand in einer sorgfältigen Revision \emph{sämmtlicher} Sterne \emph{am Himmel selbst}, bevor sie in die Karten aufgenommen wurden: freilich ein Mittel, was bei einem Unternehmen von diesem Umfange die Kräfte Eines Menschen beinahe zu übersteigen schien, aber welches nichts desto weniger Hr. Prof. \textsc{Harding} bei der Bearbeitung eines neuen Atlasses für unerlässlich hielt. So ist also in die gegenwärtigen Himmelskarten auch nicht ein einziger Stern aufgenommen worden, von dessen wirklichem Vorhandensein an seinem Platze Hr. Prof. \textsc{Harding} sich nicht vorher selbst überzeugt hätte.

Nach diesen allgemeinen Bemerkungen über den Geist, in welchem dies Unternehmen ausgeführt wird, woraus man auf dessen wissenschaftlichen Werth von selbst leicht den Schluss machen kann, wollen wir noch einige Worte über das Aeussere der Karten hinzusetzen. Hr. Prof. \textsc{Harding} hat geglaubt ein bei weitem kleineres Format, als das der \textsc{Bode}'schen Karten ist, vorziehen zu müssen, theils des bequemern Gebrauchs wegen, theils um den Preis nicht zu sehr erhöhen zu müssen: beides Rücksichten - die bei den Astronomen sehr in Betracht kommen. Die Blätter haben $16\frac{3}{4}$ Pariser Zoll Höhe, und $20\frac{1}{2}$ Zoll Breite, also nur etwa ein Drittel so viel Oberfläche, als die \textsc{Bode}'schen Karten. Bei den vier ersten Blättern, wo die Declinationen nicht über 30 Grad nördlich und südlich hinausgehen, schien es mit Recht überflüssig, eine eigentliche Projection zu wählen; die Grade der geraden Aufsteigung und der Abweichung sind also durchgehends gleich gross angenommen (etwa $5\frac{6}{7}$ Pariser Linien), welches auch für den Gebrauch noch einige Bequemlichkeiten darbietet, welche die kleine, doch selbst bei den grössten Declinationen kaum merkliche, Defiguration weit überwiegen. Jedes Blatt enthält so 42 Grad in der geraden Aufsteigung, und $34\frac{1}{2}$ Grad in der Abweichung; das erste Blatt geht in jener von $359^\circ$ bis $41^\circ$, das zweite von $39^\circ$ bis $81$ u.\,s.\,w., so dass an das neunte von $319^\circ$ bis $1^\circ$ das erste sich wieder anschliesst. Das Netz für die einzelnen Grade der Rectascension und Declination ist des bequemern Gebrauchs wegen stehen geblieben. Aber ausser diesem Netze und den Grenzen und Namen der Sternbilder enthalten die Karten nun durchaus weiter nichts, als das, was der, der sie nicht als eine Sammlung hübscher Bilder, sondern zum wirklichen Gebrauche anschafft, auf ihnen sucht, nemlich die Sterne, und diejenigen Bezeichnungen derselben, welche allgemein angenommen sind, d.\,i. die \textsc{Bayer}'schen Buchstaben oder die \textsc{Flamsteed}'schen Zahlen. Figuren der Sternbilder selbst, die bei dem heutigen Zustande der Wissenschaft ganz ausser Gebrauch sind, und also den Grund der Karten in den Augen des Kenners nur verunstaltet haben würden, sind mit Recht ganz weggeblieben, und dadurch eben die deutliche und reinliche Darstellung eines solchen Reichthums von Sternen möglich gemacht. Im Durchschnitt wird man auf jedes Blatt über 2000 Sterne zählen können, und wenn also, des Verf. Plane zufolge, das Ganze aus 24 Blättern bestehen wird (von denen der grössere Theil in der Zeichnung bereits ganz vollendet ist), so wird man in diesem Himmelsatlas ungefähr 50000 Sterne besitzen.

Vielleicht wird Mancher wünschen, dass es dem Verf. gefallen haben möchte, ausser dem Netze für Rectascension und Declination auch noch das Netz für die Längen- und Breitengrade hinzu zu fügen, \SourceLeaf{517}um besonders Planetenörter bequemer eintragen zu können, die in den astronomischen Ephemeriden gewöhnlich nur nach Länge und Breite angesetzt zu werden pflegen. Dies Netz ist indess mit gutem Vorbedacht weggeblieben, damit nemlich der Eindruck, den die Configurationen der Sterne in den Karten auf das Auge machen, so wenig als möglich durch das Gerüste schwarzer Linien gestört werde. Zum Ersatz wird aber Hr. Prof. \textsc{Harding} in dem am Ende beizufügenden Texte eine sehr geschmeidige, sorgfältig berechnete Tabelle mittheilen, wonach jeder Liebhaber jenes Netz in sein Exemplar mit der Feder sehr leicht selbst eintragen kann, und zwar nach Gefallen etwa mit einer andern Farbe, um beide Zwecke zugleich erreichen zu können. Vielleicht wäre der Verf. nicht abgeneigt, diese nützliche Tafel auch schon früher an einem schicklichen Orte für die Besitzer der ersten Lieferung mitzutheilen.

Es bleibt uns nichts übrig, als nun noch die Theile des Himmels namhaft zu machen, welche auf den vier Blättern der ersten Lieferung vorgestellt werden. Da die Vertheilung der Karten nicht nach den Sternbildern, sondern, um den Raum möglichst zu sparen, und unnöthige Wiederholungen zu vermeiden, nach den Graden der Rectascension und Declination gemacht ist, so erscheinen die meisten Sternbilder zerstückelt; dem unterrichteten Liebhaber braucht nicht erst gesagt zu werden, dass dies etwas durchaus Gleichgültiges ist. Das Blatt Nr. I., welches sich in Declination von $13^\circ$ nördl. bis $21\frac{1}{2}^\circ$ südl. erstreckt; enthält also einen Theil der Fische, des Widders, den grössten Theil des Wallfisches, und einen kleinen Theil des Eridanus; auf dem Blatt Nr. II. von $28\frac{1}{2}^\circ$ nördl. bis $6^\circ$ südl. findet man grössten Theils das Uebrige vom Widder und Wallfisch, fast den ganzen Stier und die Georgsharfe, Stücke vom Eridanus und Orion; das Blatt Nr. V. von $27^\circ$ nördl. bis $7\frac{1}{4}^\circ$ südl. gibt einen grossen Theil des Löwen, Stücke vom kleinen Löwen und Becher, den grössern Theil der Jungfrau und des Haars der Berenice; endlich liefert das Blatt Nr. IV. von $5^\circ$ nördl. bis $29\frac{1}{4}^\circ$ südl. einen Theil des Steinbocks, des Luftballons, des südlichen Fisches, der Bildhauerwerkstätte, des Wallfisches, der Fische, des Pegasus und des Füllens, und den grössten Theil des Wassermanns. -- Wir beschliessen diese Anzeige mit dem Wunsche, dass nun auch von allen, denen das Gedeihen der Wissenschaften und die Ermunterung einer echt-deutschen Beharrlichkeit in wahrhaft nützlichen Arbeiten am Herzen liegt, diese nicht genug zu empfehlende Unternehmung unterstützen und aufrecht erhalten werden, und diese zur besten Belohnung des würdigen Urhebers für die erhabene Wissenschaft künftig reichliche Früchte tragen möge.
